\documentclass[runningheads]{llncs}
%
%-------------------------------------------------------------------------------
%
\usepackage[export]{adjustbox}
\usepackage[linesnumbered,ruled,noend]{algorithm2e}
\usepackage{amsbsy}
\usepackage{amsmath}
\usepackage{amssymb}
\usepackage{dsfont}
%\usepackage{subcaption}
\usepackage{xr}
\usepackage{caption}

\usepackage{graphicx}
\usepackage{listings}
\usepackage{multicol}
\usepackage{hyperref}
\usepackage[dvipsnames]{xcolor}
\usepackage{tabularx}
\usepackage{booktabs}
\usepackage{float}
\usepackage{subfig}

\usepackage{comment}

\usepackage{multirow}

\usepackage{tikz}
\usetikzlibrary{patterns, positioning}
\usepackage{pgfplots}
\usepackage{pgfplotstable}
\pgfplotsset{compat=1.16}

\usepackage[endianness=big]{bytefield}
\usepackage{soul}
\usepackage{stmaryrd}
\pgfplotsset{every axis/.append style={
                    label style={font=\normalsize},
                    tick label style={font=\normalsize}
                    }}
\newtheorem{prop}{Property}
\pgfplotscreateplotcyclelist{alecolorlist}{%
violet \\%
blue   \\%
green  \\%
yellow \\%
orange \\%
red \\}

%%%%%%%%%%%%%
% variables used in the reference Algorithm

\newcommand{\kpub}{\ensuremath{\mathcal{K}_{pub}}}
\newcommand{\kpri}{\ensuremath{\mathcal{K}_{pri}}}
\newcommand{\sample}{\ensuremath{ \stackrel{\$}{\gets} }}
%\newcommand{\samplewt}[1]{\ensuremath{\stackrel{$}{\leftarrow}_{#1}}}
\newcommand{\ssplit}[1]{\ensuremath{ \llbracket #1 \rrbracket }}
\newcommand{\clshift}{\ensuremath{\lll_n}}


%%%%%%%%%%%%

\SetFuncSty{textsc}

%\usetikzlibrary{external}
%\tikzexternalize[prefix=tikzplots/output,mode=list and make,optimize command away=\includepdf]
%\tikzexternalize[prefix=tikzplots/output,optimize command away=\includepdf]
\renewcommand\UrlFont{\color{blue}\rmfamily}
%
%-------------------------------------------------------------------------------
%
\newcommand{\SM}[1]{\textcolor{blue}{#1}}
\newcommand{\RS}[1]{\textcolor{orange}{#1}}
\newcommand{\AB}[1]{{\color{red} #1}}
%
%-------------------------------------------------------------------------------
%
\title{A Compact and Constant Time HQC Realization on Cortex-M4}
% If the paper title is too long for the running head, you can set
% an abbreviated paper title here
\titlerunning{A Compact and Constant Time HQC Realization on Cortex-M4}
%
%-------------------------------------------------------------------------------
%
\author{Omitted for blind review}
\institute{Omitted for blind review}
%
%
%
%
%
%-------------------------------------------------------------------------------
%
\begin{document}
%
%-------------------------------------------------------------------------------
%
\maketitle              % typeset the header of the contribution
%
%-------------------------------------------------------------------------------
%
\begin{abstract}
    In this paper we show the first constant-time implementation of HQC
    targeting Cortex-M4 microcontrollers. We demostrate the efficiency
    of our implementation and its constant-time behavior empirically,
    proving that HQC is a viable candidate for the PQC standardization.\\
    We then compare the performance of our implementations with
    other post-quantum schemes, showing the efficiency of our approach.s
Key point: first Cortex-M4 optimized implementation of HQC, fully constant time
\keywords{Post-quantum cryptography, code-based cryptosystems, constant time
implementation, HQC}
\end{abstract}
%
%-------------------------------------------------------------------------------
%
\section{Introduction}\label{sec:intro}
Post-quantum cryptography (PQC), also known as quantum-resistant 
cryptography, aims at defining new public-key cryptographic schemes based 
on mathematical problems proven to be computationally hard against 
adversaries equipped with both classical and quantum computers. 
Recent advancements on quantum computer technology have witnessed a steep 
increase in scientific and engineering activities to deliver new solutions 
to some of the most pressing scientific and societal challenges~\cite{mckinsey23}. 
However, from a cybersecurity standpoint, the availability of a quantum computer 
with a sufficiently large number of qubits, would leave most of the 
current protocols used to guarantee the confidentiality and authenticity of communications and data storage vulnerable to cryptanalytic techniques based on Shor~\cite{DBLP:conf/focs/Shor94}-like quantum algorithms that subvert the 
security guarantees of the underlying public key cryptosystems.\\
\indent Since $2017$, several institutional and governmental bodies promoted the 
analysis of quantum-secure systems, but most stakeholders track the efforts 
of the PQC standardization contest promoted by the U.S. National Institute 
of Standard (NIST).  
Through two rounds of this evaluation process, NIST selected seven finalists
and eight alternate candidates, which were then thinned down to two standard
candidates and four alternates for a final round of cryptanalysis and possible 
standardization. 
The current standardization-elected candidate Key Encapsulation Method (KEM) is 
Kyber~\cite{Kyber}, a lattice-based cryptosystem, while NIST selected four 
candidates relying on different computationally hard problems, in order to obtain
resilience against potentially disruptive attacks on lattice-based cryptosystems.
Among the four candidates, a polynomial-time cryptanalytic attack eliminated
SIKE, leaving only three KEMs which are based based on the hardness of the 
syndrome decoding problem, i.e., Classic McEliece~\cite{ClassicMcEliece}, 
BIKE~\cite{BIKE} and HQC~\cite{HQC}.
Among these candidates, only HQC can exhibit a security proof 
based on a polynomial reduction of the solution of a computationally
hard problem to the breaking of the cryptosystem.
In particular, any successful break of HQC implies a faster way to solve the
syndrome decoding problem on a random quasi-cyclic code of rate $\frac{1}{2}$.
By contrast, Classic McEliece relies on an additional assumption, namely the indistinguishability of a hidden Goppa code from a random code, and BIKE relies 
on the indistinguishability of a hidden QC-MDPC code from a random quasi 
cyclic code.
While both arguments are currently unchallenged, there is no proof that they
are indeed corresponding to computationally hard problems.\\
\indent Hamming Quasi-Cyclic (HQC)~\cite{HQC,DBLP:journals/tit/Aguilar-Melchor18} is
a code-based IND-CCA secure KEM whose security relies on the syndrome 
decoding problem for quasi cyclic codes. It is obtained by applying the 
Hofheinz-H\"{o}velmanns-Kiltz (HHK) transform~\cite{DBLP:conf/tcc/HofheinzHK17} 
to a IND-CPA construction, which is also denoted as HQC.PKE.
Given the sound security record of HQC, much attention has been devoted to its 
side-channel cryptanalysis.
Most of the independent analyses focused mainly on power analysis or on 
extracting electromagnetic traces~\cite{wcc22,DBLP:conf/pqcrypto/GoyLG22,DBLP:conf/pqcrypto/SchambergerHRWS22,DBLP:conf/cardis/SchambergerRSW20};
nonetheless, in~\cite{DBLP:journals/tches/GuoHJLNS22} the authors described a 
timing-based key-recovery  attack against HQC, based on an implementation 
weakness in the step that used a procedure called \emph{rejection sampling} 
during randomness generation. 
As a countermeasure, a more careful implementation of the rejection sampling 
step is currently employed.\\
\indent 
To date, no timing and power- or Electromagnetic(EM)-based side-channel secure 
implementation of HQC, targeted to embedded systems, has been proposed, despite 
the real threat set by these attacks.
Moreover, we also note that most of the literature focused on timing side-channel
attacks aims at key retrieval, while it is also possible to perform a 
message-retrieval attack against the encryption primitive, in case the 
operations in the said encryption are not constant time.
Indeed, the attack mechanism mimics the common attacks on modular exponentiation, 
as the non-constant time polynomial multiplication techniques, such as the one
in the reference code by HQC, have a loop whose body contains branches depending
on the values of a secret operand.
Since, in the HQC.PKE encryption primitive, at least one of such operands is 
secret (as it is the plaintext message to be encrypted), providing constant time
multiplication algorithms is needed to secure the encryption primitive. 
Furthermore, as the HHK transform requires a execution of the HQC.PKE encryption 
function during the computation of the HQC KEM decryption, the requirement for 
constant time multiplication implementations is further strengthened.\\
\indent 
The current state of the software HQC implementations tailored to the ARM
Cortex-M$4$ microcontroller (chosen by NIST as the reference platform for
microcontroller benchmarks) includes only the one in PQM$4$~\cite{PQM4}, which is 
a public software project collecting the implementation of all NIST
fourth round candidates tailored for the Cortex M$4$ CPUs.
In the PQM$4$ project, the latest version of the HQC implementation is directly
derived from the third round reference one, without further optimizations.
As a consequence, the implementation is vulnerable to the timing side-channel 
attack by~\cite{DBLP:journals/tches/GuoHJLNS22}. Furthermore, no protection 
against power-based side-channel attacks is present.

\noindent\textbf{Our Contribution.} In this paper, we present the first fully 
constant-time and power/EM side-channel protected software implementation of HQC.KEM, tailored to Cortex-M$4$ microcontrollers. 
We provide an experimental validation of the security of our proposal via a 
battery of statistical $t$-tests, both for the EM-emissions and execution time 
side-channels.
We evaluate the timing side-channel immunity while varying the CPU frequency and 
memory latencies, to highlight the stability of our solution across different 
memory access latencies (which depend on the CPU frequency). 
Our results show that employing a FFT-based multiplication algorithm is the 
optimal choice for constant-time multiplications on Cortex-M$4$, in contrast 
with classic divide-et-impera techniques (e.g., Karatsuba) which provide a 
better tradeoff for x$86$\_$64$ computing platforms.
The new algorithmic strategy we introduce to hinder power/EM side-channel 
attacks against HQC.KEM software implementations achieves a lower than linear 
performance overhead in protecting the HQC.KEM Decapsulation primitive, 
currently a unique and welcome feature among side-channel protected 
implementations of post-quantum  KEMs.
Our approach is specifically tailored to the decapsulation of 
HQC, while applying an $n$ out of $n$ share-splitting strategy to 
Key-generation and Encapsulation primitives with a linear 
performance overhead.

%
%-------------------------------------------------------------------------------
%
%\section{Related Work}\label{sec:relatedWork}
%\input{related_work.tex}
%
%-------------------------------------------------------------------------------
%
\section{Background}\label{sec:background}
A Key Encapsulation Mechanism (KEM) is an asymmetric encryption technique 
designed to generate both a randomly chosen symmetric key $k_s$ and its 
encrypted version $c$, denoted as {\em encapsulated key}. 
Considering an insecure communication channel, the encapsulation $c$ serves as 
a secure delivery of the secret key, while $k_s$ is meant to be used for the 
subsequent symmetric encryption of data exchanged between the endpoints of the 
communication.
Unlike public key encryption primitives (PKEs), such as either 
RSA~\cite{DBLP:journals/cacm/RivestSA78} or the 
ElGamal~\cite{DBLP:journals/tit/Elgamal85} encryption algorithms, to prepare a proper plaintext there is no need for a criterion to map bitstrings into 
elements of an algebraic structure,
Furthermore, since the bulk of the data is encrypted with the conveyed 
symmetric key $k_s$, they can be received as a stream, without the need to 
buffer the whole ciphertext before decryption. 
From a security standpoint, the only (weak) requirement on the symmetric-key 
cipher used for subsequent data bulk encryption is that it exhibits a 
stream-based mode of operation (e.g., CTR - CounTeR mode) that is secure in case $k_s$ is employed only once per data exchange session.

Post-quantum cryptographic primitives define PKEs that are only weakly secure, 
due to the presence of probabilistic correctness errors or because they are 
proven to exhibit only the computational property of either {\em indistinguishability of ciphertexts under chosen-plaintext attack} (IND-CPA) or 
{\em one-wayness under chosen-plaintext attack} (OW-CPA).
A PKE benefits from OW-CPA if an adversary equipped with a public key and the corresponding encryption of a random plaintext, is computationally unable to 
derive the entire plaintext.
It is well known that IND-CPA implies OW-CPA, because a PKE benefits from the 
IND-CPA property when an adversary is unable to distinguish the encryptions of 
two messages of his choice; more precisely, given the public key and the 
corresponding encryption of one of two attacker-chosen plaintexts, the 
attacker is able to determine which of the two messages was encrypted only 
with negligible success probability. 
%
To guarantee security in both the Random Oracle Model (ROM) and the Quantum Random 
Oracle Model (QROM), every post-quantum cryptographic scheme makes use of a variant 
of the so-called Fujisaki-Okamoto (FO) 
transformation~\cite{DBLP:conf/crypto/FujisakiO99}, introduced for the first 
time in $1999$.
Nowadays, a newer fine-grained modular toolkit of improved FO 
transformations~\cite{DBLP:conf/tcc/HofheinzHK17} is employed for turning weakly 
secure PKE schemes into strongly secure ones, i.e., schemes benefiting from the 
property of {\em indistinguishability of ciphertexts under adaptive chosen-
ciphertext attack} (IND-CCA), i.e., an adversary is unable to distinguish the 
encryptions of two messages of his choice even if he is allowed access to a 
decryption oracle before and after seeing the (randomly selected) ciphertext of one 
of the two messages he has chosen (honoring the constraint of not querying the 
oracle with the challenge ciphertext itself).
\vspace{-0.2cm}
\subsection{HQC cryptosystem}
The first code-based PKE was introduced by McEliece in $1978$ and is based on the 
assumption that the underlying forward error-correction code, a Goppa code, used 
is indistinguishable from a random one. Although the original McEliece 
cryptosystem remains secure to this day as demonstrated by the 
highly optimized Classic-McEliece cryptoscheme~\cite{ClassicMcEliece}, the storage 
requirements for its public key, which is not smaller than $0.25$ MiB for $128$-bit 
security, is deemed too large for many application scenarios. Attempts to reduce 
the key size of the McEliece PKE by using codes more structured than Goppa codes have been all cryptanalyzed.
Therefore, the authors of the IND-CCA {\em Hamming Quasi-Cyclic} (HQC) code-based 
post-quantum KEM proposed a novel approach for the design of their PKE, which 
combines two different types of codes:
{\em (i)} An [$n, k$] code $\mathcal{C}$ defined as the concatenation of 
({\em shortened}) Reed-M\"{u}ller (RM) and ({\em duplicated}) Reed-Solomon (RS) codes, 
for which an efficient decoding algorithm is known, with a fixed, publicly known 
generator matrix $\mathbf{G} \in \mathbb{F}_2^{k \times n}$ ensuring the capability 
to correct up to $\delta$ errors (i.e., bit-flips) in a $n$-bit codeword obtained 
from a message $\mathbf{m} \in \mathbb{F}_2^k$ encoded as a (row) binary vector, as 
$\mathbf{m}\mathbf{G}$. 
{\em (ii)} A random double-circulant [$2n, n$] code with a publicly
known parity-check matrix $\mathbf{H}$$\in$$\mathbb{F}_2^{n \times 2n}$, 
$\mathbf{H}$$=$$[\mathbf{I_n} | \mathbf{rot}(\mathbf{h})]$, where $\mathbf{I_n}$ 
is the identity matrix and $\mathbf{rot}(\mathbf{h}$) is the circulant 
$n \times n$ matrix induced by the $1 \times n$  binary vector $\mathbf{h}$, 
i.e., the matrix obtained from the juxtaposition of $\mathbf{h}$ as the first 
column and the subsequent columns obtained through a vertical rotation, by one 
cell, of the previous one. 
The coefficients of the $1 \times n$  binary vector $\mathbf{h} \in \mathbb{F}_2^n$ are seen as the coefficients of a randomly chosen binary polynomial in 
$\mathcal{R}=\mathbb{F}_2[x]/(x^n-1)$, where $n$ is a prime number such that 
$\mathtt{ord}_2(n) = n-1$, i.e., $2^{n-1} = 1 \bmod n$.
%
The design rationale is that the double-circulant code is going to be used to 
add some ``noise'' (i.e., bit-flips) to the codewords of the inner RM-RS code, 
which in turn will be partially removed through the use of the secret key. The 
leftover noise plus an additional uncorrelated noise term will be decoded by the 
code $\mathcal{C}$, following the same generic steps of the ElGamal 
cryptosystem~\cite{DBLP:journals/tit/Elgamal85} and adapting to quasi-cyclic (QC) 
codes the approach first introduced in~\cite{DBLP:conf/focs/Alekhnovich03}. In such 
a way, HQC can exhibit significantly smaller keys than the Classic-McEliece KEM, 
while still providing the same computational security guarantees. Indeed, for a 
security level of $128$-bit, the public key size is around $2$ kiB vs $256$ kiB 
of the Classic-McEliece scheme.
Specifically, the PKE underlying the HQC design is defined by the following three algorithms.

\noindent \textsf{Keygen()}: samples $\mathbf{h} \xleftarrow{\mathdollar} \mathcal{R}$,  
$(\mathbf{x}, \mathbf{y}) \xleftarrow{\mathdollar} \mathcal{R}_\omega \times \mathcal{R}_\omega$, and 
outputs the private and public keypair as: \kpri = $(\mathbf{x}, \mathbf{y})$ 
and \kpub = $(\mathbf{h}, \mathbf{s} = \mathbf{H}[\mathbf{x}|\mathbf{y}]^{T} = \mathbf{x} + \mathbf{h\cdot y})$. $\mathcal{R}_\omega$ denotes the set of binary polynomials (i.e., vectors in $\mathbb{F}_2^n$) with Hamming Weight $\omega$ (i.e., the number of non-zero coefficients), while $\cdot$ is the circular convolution of the two sequences of coefficients of the factors (each of which is considered as either a binary vector or a binary polynomial); e.g., given $\mathbf{a}, \mathbf{b}, \mathbf{c}=\mathbf{a}\cdot \mathbf{b} \in \mathbb{F}_2^n$, it holds that $\mathbf{a}$$=$$(a_0, \ldots, a_i,\ldots, a_{n-1})$,  $\mathbf{b}$$=$$(b_0, \ldots, b_j,\ldots, b_{n-1})$ and $\displaystyle \mathbf{c} = \mathbf{a} \cdot \mathbf{b} = \hspace{-2.5em}\sum_{\footnotesize \begin{array}{c} k=0 \\ i+j \equiv k \bmod n\end{array}}^{n-1} \hspace{-2.5em} a_ib_j$.

\noindent \textsf{Encrypt(\kpub, \textbf{m})}: takes as input the public key, \kpub = $(\mathbf{h}, \mathbf{s} = \mathbf{x} + \mathbf{h\cdot y})$, samples an ephemeral $1 \times n$ error vector $\mathbf{e}$ from the subset of $\mathcal{R}$ that includes items with Hamming Weight $\omega_e$, i.e., $\mathbf{e} \xleftarrow{\mathdollar} \mathcal{R}_{\omega_e}$, samples the ephemeral vector $r=[\mathbf{r_1}| \mathbf{r_2}] \in \mathbb{F}_2^{2n}$ in such a way that each pair ($\mathbf{r_1}$, $\mathbf{r_2}$) is picked from the Cartesian product of two instances of the subset of $\mathcal{R}$ that includes items with Hamming Weight  $\omega_r$, i.e., $(\mathbf{r_1}, \mathbf{r_2}) \xleftarrow{\mathdollar} \mathcal{R}_{\omega_r} \times \mathcal{R}_{\omega_r}$, and output a ciphertext $\mathbf{c}$ which, in turn, is composed by the pair $(\mathbf{u}, \mathbf{v})$, where: 
$$\mathbf{u} = \mathbf{r_1} + \mathbf{h}\cdot\mathbf{r_2},\ \text{and}\ \mathbf{v} = \mathbf{mG} + \mathbf{s\cdot r_2} + \mathbf{e}$$

\noindent \textsf{Decrypt(\kpri, $(\mathbf{u}, \mathbf{v})$)}: takes as input the private key, \kpri = $(\mathbf{x}, \mathbf{y})$ and a ciphertext $(\mathbf{u}, \mathbf{v}) \in \mathbb{F}_2^n \times \mathbb{F}_2^n$, and outputs the original plaintext message through computing $\mathbf{m} = \mathcal{C}.\textsf{Decode}(\mathbf{v} - \mathbf{u\cdot y})$.

\smallskip
\noindent The computational security of recovering either the secret-key or the plaintext of an instance of the HQC.PKE scheme is computationally equivalent to solving the syndrome decoding problem on random quasi-cyclic codes with length $n$ and dimension $k$~\cite{DBLP:journals/tit/Aguilar-Melchor18}. 
While random quasi-cyclic codes are not proven to yield an NP-complete and hard 
to approximate syndrome decoding problems, it is largely believed that, in case 
$n$ is chosen as a prime such that $\mathtt{ord}(2)$$=$$n$$-$$1$, no speedup 
can be achieved with respect to the solution of the syndrome decoding problem
on random codes with the same length, when the error generating the syndrome 
is half of the code distance expected for a random code~\cite{DBLP:journals/tit/Aguilar-Melchor18}.

Thus, the best classical and quantum cryptanalytic techniques that are currently 
known to perform a key- or a message-recovery attack against the cryptosystem 
have been considered to design the parameters sets shown in Table~\ref{tab:params}. 
In particular, the NIST call defines the security levels for which parameter
sets should be designed stating that the resulting computational effort to 
break the cryptoscheme should match the one to break both classical and quantum realizations of AES-$128$, AES-$192$, and AES-$256$, respectively.
The sets of parameters of a post-quantum cryptosystem that match the security 
desiderata are labeled by NIST as Category $1$, Category $3$ and Category $5$, 
respectively, while the authors of HQC.KEM named the corresponding instances of 
their proposal as \texttt{hqc-128}, \texttt{hqc-192} and \texttt{hqc-256}.
% 
Finally, the most up-to-date specification of HQC reports a public code that is 
defined as the concatenation of Reed-Solomon (outer) and Reed-M\"{u}ller (inner) 
codes, with codeword length equal to $n_1$ and $n_2$, respectively, while the overall 
codeword length $n$ is taken as the smallest prime number greater than $n_1n_2$ 
such that $\mathtt{ord}_n(2)$$=$$n$$-$$1$. Thus, the public generator matrix 
$\mathbf{G}$ has dimension $k\times n$, where $k$ coincides with the bitsize of 
the plaintext message $\mathbf{m}$. Furthermore, the random code with 
parity-check matrix $\mathbf{H}$ is instantiated by randomly sampling the 
vector $\mathbf{h}$ from $\mathds{F}_2^n$. 
\begin{table}[!t]
    \centering
    \caption{HQC Parameter sets. $n$: public code length, $n_1$: Reed–Solomon code length, $n_2$: Reed–M\"{u}ller code length. $\omega$: weight of the $n$-bit binary vectors $\mathbf{x}, \mathbf{y}$, composing the secret key. 
    $\omega_r=\omega_e$: weight of the $n$-bit binary vectors used in the 
    encryption. $[n, k, d]$ for Reed–Solomon and Reed–M\"{u}ller codes are the 
    length, dimension, and minimum distance of the codes, respectively.
    The multiplicity of duplicated Reed–M\"{u}ller code is $3$, $5$, $5$ for \textsc{hqc-128}, \textsc{hqc-192}, \textsc{hqc-256}.
    The probability of a decapsulation error (i.e., of a correctness error of the system) is specified as the decoding failure rate (DFR) of the public code}
    {\small
     \begin{tabular}{cc|cccc|cc|c}
        \toprule
        \textbf{NIST} & \multirow{2}{*}{\textbf{Instance}} & \textbf{Plaintex} & \multirow{1}{*}{\bf 
        Reed-} & \multirow{1}{*}{\bf Reed-} & \multirow{2}{*}{$\mathbf{n}$} & \multirow{2}{*}{$\mathbf{\omega}$} & \multirow{2}{*}{$\mathbf{\omega_r = \omega_e}$} & \multirow{2}{*}{\textbf{DFR}}\\
        \textbf{Category} &  & \textbf{size (bit)} & {\bf Solomon} & {\bf M\"{u}ller}&  &  &  & \\
        \midrule
        $1$& \texttt{hqc-$128$} & $128$ &$[46, 16, 15]$ &  $[384, 8, 192]$ & $17,669$ & $66 $ & $ 75$ & $2^{-128}$\\
        $3$& \texttt{hqc-$192$} & $192$ &$[56, 24, 16]$ &  $[640, 8, 320]$ & $35,851$ & $100$ & $114$ & $2^{-192}$\\
        $5$& \texttt{hqc-$256$} & $256$ &$[90, 32, 29]$ &  $[640, 8, 320]$ & $57,637$ & $131$ & $149$ & $2^{-256}$\\
        \bottomrule
    \end{tabular}
    \label{tab:params}
    }\vspace{-0.4cm}
\end{table}

%
%-------------------------------------------------------------------------------
%
\section{Constant time, High Efficiency HQC}\label{sec:implem}
The HQC.KEM is obtained starting from the HQC.PKE, augmenting it with components
which do not involve polynomial arithmetic operations. 
In particular, the HQC.KEM employs the \textsc{SHAKE} extendable output function~\cite{SHAandSHAKE} (i.e., a function mapping an arbitrary-length bit 
string to a bitstring with a user chosen length, which is indistinguishable from
a randomly selected string; in other words a PRNG algorithm with a seed of 
arbitrary length) to realize the variant of the FO 
transform~\cite{DBLP:conf/tcc/HofheinzHK17}. Since the constant time and 
side channel protection of \textsc{SHAKE} have been successfully tackled in the 
literature, we consider them out of scope for this work, 
and employ the existing results.
Except for \textsc{SHAKE}, the HQC.KEM algorithmic components are the same of the 
HQC.PKE, namely polynomial addition and multiplication in $\mathcal{R}$, 
pseudorandom bit generation via \textsc{SHAKE} and encoding/decoding for the 
concatenated Reed-M\"{u}ller Reed-Solomon code.
Considering the fact that binary polynomial additions are straightforward to 
compute in constant time, and protect with respect to side channel attacks, 
we focus on polynomial multiplications and encoding/decoding. The pseudorandom
number generation component is effectively and efficiently protected in the
same way the HQC.KEM \textsc{SHAKE} element is.

Concerning the encoding/decoding, we devise an algorithmic strategy to avoid
the requirement for a power based side channel attack protected decoder, 
exploiting the linearity of the RM-RS codes, while no timing analysis concerns 
are present thanks to their constant time implementation being the current one 
at the state of the art.
In the following, we describe in detail the approaches to a timing-side channel
protected polynomial multiplication, while deferring the protection of the said
multiplication against power analysis attacks to the description of our 
algorithmic strategy for the protection of the entire HQC.PKE.

\subsection{Polynomial Multiplication Approaches}
HQC.PKE uses polynomial multiplications in $\mathcal{R}=\mathbb{F}_2[x] / (x^n-1)$
in each of its primitives. Referring to the definition of the HQC.PKE 
primitives and to Table~\ref{tab:params}, \textsf{Keygen()} requires one 
multiplication between elements in $\mathcal{R}$, where the polynomial 
employed as a second factor has only $\omega$ of non-null coefficients, 
with $\omega \ll n$, i.e., $\mathbf{h} \cdot \mathbf{y}$, with 
$\mathbf{h} \in \mathcal{R}$, $\mathbf{y} \in \mathcal{R}_{\omega}$. 
Thus, in any software implementation, the first factor (e.g., $\mathbf{h}$) 
is stored as a {\em dense} polynomial, i.e., a sequence of $n$ binary 
coefficients, while the second one ($\mathbf{y}$) is stored as a {\em sparse} 
polynomial, i.e., as a sequence of integer numbers, each of which is intended 
as the exponent of a non-null monomial ($x^i, 0 \leq i \leq n-1$). 
\textsf{Encrypt()} requires two multiplications in $\mathcal{R}$ to compute the 
ciphertext, in each of which the first factor is a dense polynomial and the 
second factor is a sparse polynomial with a number of coefficients equal to 
$\omega_r$. 
Finally, \textsf{Decrypt()} requires one multiplication in $\mathcal{R}$, where 
the first factor is a dense polynomial ($\mathbf{u}$) and the second one is a 
sparse polynomial ($\mathbf{y}$) with $\omega$ non-null coefficients.\\ 
Denoting as $\mathtt{wsize}$ the bitlength of the word of the CPU employed as a 
target computational platform, the data structure for storing a dense polynomial 
$b(x) \in \mathcal{R}$ is an array $B = (B_{t-1}, \ldots, B_j, \ldots, B_0)_{\mathtt{w}}$ of $t = \left\lceil \frac{n}{\mathtt{wsize}} \right\rceil$ cells, where each cell 
$B_j$ stores a sequence of adjacent binary coefficients in $b(x)$ that fits in 
a CPU word. Equivalently, the same polynomial can be thought of as a 
multiple precision unsigned integer encoded in $t$ CPU words, with each digit 
(limb) $0$$\leq$$B_j$$<$$\mathtt{w}$, $0$$\leq$$j$$\leq$$t$$-$$1$, meant to 
follow the usual positional notation of a number in radix 
$\mathtt{w} = 2^{\mathtt{wsize}}$.
A sparse polynomial $a(x) \in \mathcal{R}$ is an array $A = (A_{\omega-1}, \ldots, A_i, \ldots, A_{0})$, with $\omega \ll n$ cells, such that each cell $A_i=a$ is an 
unsigned integer that encodes the exponent of a non-null monomial 
$x^a \in \mathcal{R}$ included in $a(x)$, where the value $0\leq a\leq n-1$ is 
encoded with at most $m=\lceil \log_2(n)\rceil$ bits. 
The product $c(x) = b(x) \cdot x^{A_i} \bmod (x^n-1)$ between a dense polynomial 
$b(x)$ and a monomial with $A_i = a$  must be equivalent to the computation of 
a circular left-shift of the coefficients of $b(x)$, for $a$ positions across 
their original $n$-bit positions (also denoted as $\clshift a$). 

\subsection{Sparse-to-dense Multiplication Algorithm}\label{subsec:shift_add}
% \begin{figure*}[!t]
% \noindent\begin{minipage}{0.99\textwidth}
% \footnotesize
% \begin{algorithm}[H]
%     \caption{Sparse-to-dense Polynomial Multiplication\label{alg:varTimeShiftAndAdd}}
%     \KwIn{$n$: maximum number of bits for a polynomial in $\mathcal{R}$$=$$\mathbb{F}_2[x]/(x^n$$-$$1)$.\newline 
%     $\mathtt{wsize}$: bitlength of a CPU word, $\mathtt{w} = 2^{\mathtt{wsize}}$ radix employed to represent in positional notation a long integer/binary-polynomial.\newline   
%     $A = (A_{\omega-1}, \ldots, A_i, \ldots, A_{0})$, represents $a(x) \in \mathcal{R}$, where $A_i$ is an unsigned integer that encodes the exponent of a non-null monomial $x^a \in \mathcal{R}$, $0\leq a\leq n-1$, $\omega < n$.\newline 
%           $B = (B_{t-1}, \ldots, B_j, \ldots, B_0)_{\mathtt{w}}$, $t = \left\lceil \frac{n}{\mathtt{wsize}} \right\rceil$, represents $b(x) \in \mathcal{R}$, where a limb $B_j$ stores a sequence of binary coefficients fitting in a CPU word having $\mathtt{wsize}$ bits (i.e., $0$$\leq$$B_j$$<$$\mathtt{w}$), $0$$\leq$$j$$\leq$$t$$-$$1$.} 
%     \KwOut{$C = (C_{t-1}, \ldots, C_k, \ldots, C_0)_{\mathtt{w}}$, represents $c(x) = a(x)\cdot b(x) \in \mathcal{R}$, where each limb $C_k$ stores a sequence of binary coefficients fitting into a CPU word (i.e., $0$$\leq$$C_k$$<$$\mathtt{w}$), $0$$\leq$$k$$\leq$$t$$-$$1$.}
%     \BlankLine
%     $C \gets (0,\ldots,0)$ \tcp*[l]{\color{blue}zeroed array with $t$ cells, each with\ \text{{\em wsize}}\ bits}
%     \For{$i \gets 0$ \KwTo $\omega-1$}{
%        $\mathtt{tmp} \gets (\ B \clshift A_i\ )$\tcp*[l]{\color{blue}\footnotesize variable-time comput. of a $t$-word array}    
%        \For{$k \gets 0$ \KwTo $(t-1)$} {
%            $C_k \gets C_k \oplus \mathtt{tmp}_k$\;
%         }
%     }
%     \BlankLine
%     \KwRet $C$\;
%     \end{algorithm}
% \end{minipage}
% \vspace{-0.4cm}
% \end{figure*}
In the third round of the NIST standardization process, the authors of HQC 
proposed, as a reference implementation, only a schoolbook non-constant time 
polynomial multiplication algorithm using a minor variant of the schoolbook
\textit{Shift-and-Add} method, exploiting the sparsity of one of the two 
operands.
The approach is summarized as $c(x) = b(x) \cdot a(x) = b(x) (x^{A_{\omega-1}}\oplus\cdots \oplus x^{A_{i}}\oplus\ldots)$ $= b(x)x^{A_{\omega-1}} \oplus \cdots\oplus b(x)x^{A_{i}} \cdots\oplus b(x)x^{A_{0}}$, where $A_{\omega-1},\ldots, A_{0} \in \{0,\ldots ,n-1\}$ 
and each multiplication by a monomial is computed as a circular left-shift of 
the coefficients of the first polynomial $b(x)$ for $A_i$ positions across its 
original $n$ bit positions. 
The pseudo-code of the method is in Algorithm~\ref{alg:varTimeShiftAndAdd}, Appendix~\ref{sec:appendix}, while its worst-case computational 
complexity is ${O}(\omega(2t)$) bit operations.

The major drawback of the schoolbook \textit{Shift-and-Add} method lies in the 
variable-time computation reported on Alg.~\ref{alg:varTimeShiftAndAdd}, 
$\mathbf{line}\ 3$, which reveals on the timing side-channel, the sensitive 
value of each monomial coefficient $A_i$ of the sparse polynomial factor. 
Indeed, the computation time of a circular left-shift for $a$ positions across 
the original $n$ bit positions of the first array is a multiple of the execution 
time required to compute ($B \clshift 1$).         
% 
\begin{figure*}[!t]
\noindent\begin{minipage}{0.99\textwidth}
\footnotesize
\begin{algorithm}[H]
\SetAlgoLined
    \caption{Constant-time Sparse-to-dense Polynom. Multiplication\label{alg:CTShiftAndAdd}}
    \KwIn{$n$: maximum number of bits for a polynomial in $\mathcal{R}$$=$$\mathbb{F}_2[x]/(x^n$$-$$1)$.\newline 
    $\mathtt{wsize}$: bitlength of a CPU word, $\mathtt{w} = 2^{\mathtt{wsize}}$ radix employed to represent in positional notation a long integer or the sequence of coefficients of a binary-polynomial.\newline   
    $A = (A_{\omega-1}, \ldots, A_i, \ldots, A_{0})$, represents $a(x) \in \mathcal{R}$, where $A_i=a$ is an unsigned integer that encodes the exponent of a non-null monomial $x^a \in \mathcal{R}$, $0\leq a\leq n-1$, $\omega < n$, $a = (a_{m-1}, \ldots, a_0)_{\mathtt{bin}}$, $a_i \in \{0,1\}$, $m = \lceil \log_2(n) \rceil$.\newline 
    $B = (B_{t-1}, \ldots, B_j, \ldots, B_0)_{\mathtt{w}}$, $t = \left\lceil \frac{n}{\mathtt{wsize}} \right\rceil$, represents $b(x) \in \mathcal{R}$, where a limb $B_j$ stores a sequence of binary coefficients fitting in a CPU word having $\mathtt{wsize}$ bits (i.e., $0$$\leq$$B_j$$<$$\mathtt{w}$), $0$$\leq$$j$$\leq$$t$$-$$1$.} 
    \KwOut{$C = (C_{t-1}, \ldots, C_k, \ldots, C_0)_{\mathtt{w}}$, represents $c(x) = x^a \cdot b(x) \in \mathcal{R}$, where each limb $C_k$ stores a sequence of binary coefficients fitting into a CPU word (i.e., $0$$\leq$$C_k$$<$$\mathtt{w}$), $0$$\leq$$k$$\leq$$t$$-$$1$.}
    \BlankLine
    \For{$i \gets 0$ \KwTo $\omega-1$}{
    $a \gets A_i$\tcp*[l]{\color{blue}value with $m$$=$$\lceil \log_2(n)\rceil$ bits: $(a_{m-1},\ldots,m_0)_{\mathtt{bin}}$}
    $C \gets (\ B \clshift (a \mod \mathtt{wsize}) \ )$\tcp*[l]{\color{blue}\footnotesize computation of a $t$-word array}
    \For{$j \gets \log_2(\mathtt{wsize})$ \KwTo $(m-1)$} {
        $\mathtt{tmp} \gets (\ B \clshift 2^j\ )$\tcp*[l]{\color{blue}\footnotesize computation of a $t$-word array}
        $\mathtt{mask} \gets 0 - (a_j = 1)$\tcp*[l]{\color{blue}\footnotesize single word: all its bits are $1$ or $0$}
        \For{$k \gets 0$ \KwTo $(t-1)$} {
           $C_k \gets C_k \oplus ((C_k \wedge \mathtt{mask}) \vee (\mathtt{tmp}_k \wedge \mathtt{mask}))$\;
        }
    }
    }
    \KwRet $C$\;
    \end{algorithm}
\end{minipage}
\vspace{-0.4cm}
\end{figure*}
A constant-time realization of the Shift-and-Add strategy is shown in Algorithm~\ref{alg:CTShiftAndAdd}. The rationale is still to split the computation 
as the iterated sum of the dense polynomial multiplied by a monomial one. 
However, instead of executing in variable-time every single monomial 
multiplication, it is possible to realize it as a fixed sequence of circular 
left-shift operations followed by arithmetic predicates for the conditional 
committment of each of them.
Specifically, considering $a$$=$$(a_{m-1}, \ldots, a_j, \ldots, a_0)_{\mathtt{bin}}$, $0$$\leq$$j$$<$$m$, $a_j$$\in$$\{0,1\}$, $m$$=$$\lceil \log_2(n) \rceil$, 
the multiplication $c(x) = b(x) \cdot x^a \bmod (x^n$$-$$1)$ is decomposed as: 
\begin{center}
$c(x) = b(x)  \cdot \left( x^{a_{m-1}2^{m-1}} \cdots x^{a_{j}2^{m-1}} \cdots x^{a_{0}} \right)\bmod (x^n-1) $ $=\phantom{MMMMMMM...}$\newline
$= b(x)  \cdot \left( x^{a_{m-1}2^{m-1}} \cdots x^{a_{\log_2(\mathtt{wsize})}2^{\log_2(\mathtt{wsize})}} \cdots x^{a \bmod \mathtt{wsize}} \right) \bmod (x^n-1)$, 
\end{center}
observing that $\mathtt{wsize}$ is a power of two (for any CPU), thus $a = (a_{m-1},  \ldots, a_0)_{\mathtt{bin}}$ $= (a_{m-1}, \ldots, a_{\log_2(\mathtt{wsize})})_{\mathtt{bin}}$ concatenated with the binary string encoding $(a \bmod \mathtt{wsize})$, while the execution of $b(x) \cdot x^{a \bmod \mathtt{wsize}}$ is amenable to be realized through a constant amount of $t=\lceil \frac{n}{\mathtt{wsize}}\rceil$ circular left-shifts. These observations, together with  the arithmetic test to establish whether any $a_j$ is asserted or not, and a conditional update of the final result (see Alg.~\ref{alg:CTShiftAndAdd}: \textbf{lines} $4$--$8$), allow to realize a constant-time version of the Shift-and-Add method with a computational cost that amounts to $\Theta\left(\omega\,(2t)\,m + \omega\, t+\omega\, m\right)$ bit operations. 

\subsection{Dense-to-dense Multiplication Algorithms}\label{subsec:CTDenseMul}
Another approach to constant-time polynomial multiplication is to give up the 
savings offered by the sparse representation of one of the factors and represent
both factors as dense polynomials. 
Indeed, the execution of algorithmic strategies as the schoolbook multiplication 
(recommended for operands encoded with up to $\approx 300$ bits), the 
Karatsuba-, Toom-Cook- or Fast-Fourier-Transform-based (FFT) multiplications 
(recommended for operands with bitlength up to $\approx 10^4$, from 
$\approx 10^4$ to $\approx 6\cdot10^4$, and $\gtrapprox 6\cdot10^4$, 
respectively) exhibit no information leakage on the timing side-channel. 
 
\noindent
\textbf{HQC Approach.}
% From a performance perspective, the choice of the algorithmic strategy that better fits the multiplication operations listed by the HQC primitives depends on the bitlength of the operands, which in turn depend on the HQC instance that is specifically considered among \texttt{hqc-128}, \texttt{hqc-192}, and \texttt{hqc-256}.
In the fourth round of the NIST standardization process, given the size of the 
polynomials listed in Table~\ref{tab:params}, for the realization of polynomial 
multiplications required in \texttt{hqc-128}, \texttt{hqc-192}, the authors of 
HQC decided to apply a combination of the Toom-Cook multiplication strategy 
named as Toom-$3$, which splits each factor in three parts before performing 
operations on them, with a recursive $3$-split Karatsuba multiplication approach 
or a recursive $5$-split Karatsuba multiplication approach, down to the point 
where a $64$-bit$ \times 64$-bit multiplier is sufficient. 
For the realization of polynomial multiplications required in \texttt{hqc-256} 
the same overall approach is applied, replacing the $3$-split Karatsuba 
algorithm with a $5$-split one. The chosen strategy is supported by a favorable 
$O(n^{1.46})$ cost of the Toom-Cook approach, which is further optimized, in 
the finite regime of the HQC parameters, applying also the Karatsuba technique 
and the (vectorized) multiply instructions provided by the instruction set 
architecture of the target platform.
 
\noindent
\textbf{Alternate Approach.}
The asymptotically fastest methods currently known to perform multiplications in 
$\mathcal{R}$ require to evaluate them at $2n$ points, perform a pointwise 
multiplication of the these results, and interpolate the resulting $2n$ values 
to get the final polynomial. While this approach may seem less efficient, 
with a proper choice of the evaluation points it takes the name of 
Fast Fourier Transform (FFT) based multiplication, and the overall 
computational cost (including the one of interpolation) requires 
${O}(n\cdot\log(n))$ bit operations.

We noted that relatively recent development in the field of FFT-based arithmetic 
in binary finite fields may allow execution of polynomial multiplications in 
both $\mathtt{F}_2[x]$ and $\mathcal{R}$ exhibiting a computational complexity 
that is less than the one that can be obtained with other methods when the 
bitlength of the operands in in the range prescribed by each HQC instance.
% 
Specifically, Cantor~\cite{DBLP:journals/jct/Cantor89}, developed the so-called 
additive-FFTs through evaluating polynomials at points that form an additive 
subgroup in $\mathbb{F}_2$ and defining a basis, termed {\em Cantor basis}, for 
constructing an FFT over the field that works with the complexity of 
$O(n \log(n))$ multiplications and 
$O(n\cdot\log^{\log_2 (3)}(n))$ $\mathtt{xor}$-additions for evaluating 
$n = 2^m$ elements. 
In $2010$, Gao and Mateer~\cite{DBLP:journals/tit/GaoM10} presented an 
additive-FFT (GM-FFT) over $\mathbb{F}_{2^k}$ with 
$O(n\cdot\log^2(n))$ $\mathtt{xor}$-operations for evaluating a polynomial at 
$n = 2^m $ points in general. 
It can be optimized to $O(n\cdot\log( n)\cdot\log (\log(n)))$ only when $m$ is 
a power of $2$, and the polynomials in this case are represented in a special 
basis introduced by Cantor~\cite{DBLP:journals/jct/Cantor89}. 
In $2014$, Bernstein and Chou~\cite{DBLP:conf/sacrypt/BernsteinC14} presented an 
efficient implementation of the GM-FFT with performance equal to the ones 
obtainable applying the more general variant introduced in the same year by 
Lin, Chung, and Han~\cite{DBLP:conf/focs/LinCH14} (named LCH-FFT ), which uses 
a polynomial basis different from the standard one. 
In the subsequent work~\cite{DBLP:journals/tit/LinAH16} a method to convert 
from/to the standard polynomial basis to/from the one required 
by~\cite{DBLP:journals/jct/Cantor89} is introduced. 
The computational complexity to realize the conversion is 
$O(n\cdot \log(n)\cdot \log(\log(n)))$ $\mathtt{xor}$-operations, with $n$ 
being any power of two, which in turn reaches the same complexity as the 
common multiplicative FFTs.

We were able to effectively reproduce the multiplication strategy introduced  
in~\cite{DBLP:conf/sacrypt/BernsteinC14} to the end of realizing the polynomial 
multiplications listed in the HQC primitives with a constant-time dense-to-dense
algorithmic approach.
Since the length of the operands $n$ in HQC is not a power of two, a 
straightforward approach would call for zero-padding of the polynomials 
for a number of coefficients up to the nearest power of two greater than $n$.
However, this approach performs a significant amount of unnecessary computations
among zero-filled words, whenever $n$ is far from the said power of two (as it
is the case in \textsc{hqc-128} and \textsc{hqc-192}).
To avoid such overhead, we automatically generated the code of the  
multiplication algorithm in~\cite{DBLP:conf/sacrypt/BernsteinC14}, removing at code
generation time the operations acting deterministically only on zero-filled words.\footnote{The source code of our HQC realization will be publicly available upon 
acceptance} 

\subsection{Boolean Masking of the HQC.PKE against Power Analysis}\label{subsec:masking}
Regardless of their security against mathematical cryptanalysis, the 
implementation of cryptographic algorithms on end-user devices may leak 
information about secret parameters of the scheme, which can lead to 
successful plaintext or secret-key recovery attacks. 
Gathering power- or EM-measurements (a.k.a. {\em traces}) during a known-input 
execution of a cryptographic primitive allows to apply a divide-and-conquer 
strategy in inferring the secret values in the computation, thanks to the 
timewise locality and granularity of the information obtained.
The general workflow for non-profiled side channel attacks predicts the side
channel behaviour of a simple operation in the cryptographic primitive, depending
on a (guessed) value for the small portion of secret involved in it.
After collecting measurements from the device, the correct guess for the secret
is singled out comparing with statistical means the secret-dependent predictions
with the actual measures.

A popular, and quantifiably secure countermeasure against side-channel attacks 
is \emph{masking}~\cite{ishai2003private,prouff2013masking}.
Masking works through performing the entire computation on share-split values, 
through an $x$-out-of-$x$ share splitting schemes.
The number of shares into which a value is split provides a quantification
of the security margin: this is commonly denoted as the \emph{masking order} $d$, 
where a masking scheme with $x$ shares is known as a $d=x-1$ order masking.
The most common share splitting strategy is the one known as 
\emph{Boolean masking}.
Applying Boolean masking to a polynomial in $\mathbf{v}\in\mathcal{R}$ we have that the $x$-uple $\ssplit{\mathbf{v}}= (\mathbf{v}_0,\ldots,\mathbf{v}_{x-1})$ of the shares of a polynomial $\mathbf{v}\in \mathcal{R}$ is obtained drawing $x-1$ random polynomials, 
$(\mathtt{rnd}_0,\ldots,\mathtt{rnd}_{x-2})$ and obtaining the remaining share 
as their coefficient-wise \texttt{xor} sum with the polynomial to be 
split $\mathbf{v}$, $(\mathbf{v}\oplus \bigoplus_{i=0}^{x-2} 
\mathtt{rnd}_i)$. We denote this operation as 
$\text{\sc Split}(\mathbf{v}) = (\mathbf{v}\oplus \sum_{i=0}^{x-2} \mathtt{rnd}_i,\mathtt{rnd}_0,\ldots,\mathtt{rnd}_{x-2})=\ssplit{\mathbf{v}}$.

All the operations on share-split polynomials must be performed so that combining
the share-split results (i.e., \texttt{xor}-ing together all shares) yields the correct result. 
We will denote such an action on the share-split value $\ssplit{\mathbf{v}}$ as
$\text{\sc Combine}(\ssplit{\mathbf{v}})=\mathbf{v}$.
The share combination is performed only at the end of the primitive computation, as any operation combining all the shares from of a single original value $\mathbf{v}$ reveals information on the value itself.
By contrast, combining up to $x-1$ shares does not yield any information 
leakage.

Boolean masking incurs in only a linear (in $x$) overhead whenever the performed 
operations on share-split polynomials are \texttt{xor} linear, nonlinear operations (e.g., Boolean \texttt{and}) have a quadratic overhead~\cite{ishai2003private}.

\begin{figure*}[!t]
\noindent\begin{minipage}{0.45\textwidth}
\footnotesize
    \begin{algorithm}[H]
        \caption{Masked HQC.PKE \textsf{Keygen}\label{alg:keygen}}
        \KwIn{none}
        \KwOut{$\kpub = (\mathbf{h},\mathbf{s})\in \mathcal{R}\times \mathcal{R}$\newline 
        $\kpri$$=$$(\mathbf{x},\mathbf{y})$$\in$$\mathcal{R}_\omega$$\times$$\mathcal{R}_\omega$}

        $\mathbf{h} \sample \mathcal{R}$\; 
        $(\mathbf{x},\mathbf{y}) \sample \mathcal{R}_\omega \times \mathcal{R}_\omega$\;
        $\ssplit{\mathbf{y}} \gets \text{\sc Split}(\mathbf{y})$\;
        $\ssplit{\mathbf{s}} \gets (\mathbf{h}\cdot\mathbf{y}_{,0},\ldots,\mathbf{h}\cdot\mathbf{y}_{,x-1})$\;
        $\ssplit{\mathbf{x}} \gets \text{\sc Split}(\mathbf{x})$\;
        $\ssplit{\mathbf{s}} \gets \ssplit{\mathbf{x}} \oplus \ssplit{\mathbf{s}}$  \;
        $\kpri \gets (\mathbf{x},\mathbf{y})$\;
        $\kpub \gets (\mathbf{h},\textsc{Combine}(\ssplit{\mathbf{s}}))$\;
        \KwRet $\kpri$, $\kpub$\;
        \BlankLine
        \BlankLine
        \BlankLine
        \BlankLine
        \BlankLine
        \BlankLine
        \BlankLine
        \BlankLine
        \BlankLine
        \BlankLine
        \BlankLine
    \end{algorithm}
\end{minipage}\hspace{0.5cm}
\begin{minipage}{0.5\textwidth}
\footnotesize
    \begin{algorithm}[H]
        \caption{Masked HQC.PKE \textsf{Encrypt}\label{alg:enc}}
        \KwData{$\mathbf{G}$: a $\mathds{F}_2^{k \times n}$ matrix generating an efficiently correctable code}
        \KwIn{$\kpub = (\mathbf{h},\mathbf{s})\in \mathcal{R}\times \mathcal{R}$\newline 
        $\mathbf{m}\in \mathbb{F}_2^k$}
        \KwOut{$(\mathbf{u}, \mathbf{v}) \in \mathcal{R} \times \mathcal{R}$}

        $\mathbf{e} \sample \mathcal{R}_{\omega_e}$, 
        $(\mathbf{r_1},\mathbf{r_2}) \sample \mathcal{R}_{\omega_r} \times \mathcal{R}_{\omega_r}$\;
        $\ssplit{\mathbf{r_1}} \gets \text{\sc Split}(\mathbf{r_1})$\;
        $\ssplit{\mathbf{r_2}} \gets \text{\sc Split}(\mathbf{r_2})$\;
        $\ssplit{\mathbf{u}} \gets \ssplit{\mathbf{r_1}} \oplus 
        (\mathbf{h}\cdot\mathbf{r_2}_{,0},\ldots,\mathbf{h}\cdot\mathbf{r_2}_{,x-1})$\;
        $\ssplit{\mathbf{m}} \gets \text{\sc Split}(\mathbf{m})$\;
        $\ssplit{\mathbf{v}} \gets \ssplit{\mathbf{m}}\mathbf{G}  \oplus \text{\sc Split}(\mathbf{e}) $\;
        $\ssplit{\mathbf{v}} \gets \ssplit{\mathbf{v}} \oplus (\mathbf{s}\cdot\mathbf{r_2}_{,0},\ldots,\mathbf{s}\cdot\mathbf{r_2}_{,x-1}) $\;
        $\mathbf{u} \gets \textsc{Combine}(\ssplit{\mathbf{u}})$\;
        $\mathbf{v} \gets \textsc{Combine}(\ssplit{\mathbf{v}})$\;
        \KwRet $(\mathbf{u},\mathbf{v})$\;
    \end{algorithm}
\end{minipage}
\vspace{-0.4cm}
\end{figure*}
To the end of devising our algorithmic proposal for efficient Boolean masking
of HQC.PKE, we observe that the following two properties hold for its structure:
\begin{prop}[Public operands]\label{prop:one} All the polynomial multiplications have a factor 
which is public and either input independent or constant (across a computation). 
It is thus possible to avoid the share-splitting of the public operand and 
rewrite the multiplication as a sequence of $x$ share-wise \texttt{xor} 
additions of the share-split operand and $x$ share-wise circular shifts.
\end{prop}
\begin{prop}[Share-split decoding]\label{prop:two} Computing in a share-split 
representation the decoding w.r.t. the public code $\mathbf{G}$ is  
performed in a more efficient way exploiting the linearity of the encoding: 
the $x-1$ random shares for the erroneous codeword to be decoded are obtained 
encoding $x-1$ random information words.
\end{prop}
\noindent\textbf{Masked} \textsf{Keygen}.
The masking of the key generation algorithm (Algorithm~\ref{alg:keygen}) is 
performed splitting the private key ($\kpri = (\mathbf{x},\mathbf{y}) \in 
\mathcal{R}_\omega \times \mathcal{R}_\omega$) element $\mathbf{v}$ (line $3$), 
and multiplying it by the public polynomial $\mathbf{h}$, which does not need
to be share-split (Property~\ref{prop:one}). To add the secret value 
$\mathbf{x}$ to $\ssplit{\mathbf{s}}$, we share-split it (line $5$) in turn 
and perform a share-wise \texttt{xor}. Finally, we combine the shares of 
$\ssplit{\mathbf{s}}$ to obtain the second half of the public key.

\noindent\textbf{Masked} \textsf{Encrypt}. In the encryption algorithm, 
we protect both the message $\mathbf{m}$ and the ephemeral values 
$\mathbf{r_1},\mathbf{r_2},\mathbf{e}$ by share splitting them (lines $2,3,5,6$).
The multiplications at lines $4$ and $6$ can be performed without share splitting
the public factors $\mathbf{h}$ and $\mathbf{s}$ -- by Property~\ref{prop:one}, 
indeed, they are the two components of the public key -- and therefore their 
share-wise multiplication only takes linear time in the number of shares $x$.
The encoding of the shared message $\ssplit{\mathbf{m}}$ is performed via a sequence
of linear operations, i.e. the share-wise multiplication by the public binary 
matrix $\mathbf{G}$, which therefore does not need to be share split, or any 
equivalent encoding technique.
\begin{figure*}[!t]
\noindent\begin{minipage}{0.99\textwidth}
\footnotesize
\begin{algorithm}[H]
    \caption{Masked HQC.PKE \textsf{Decrypt}\label{alg:dec}}
    \KwData{$\mathbf{G}$: a $\mathbb{F}_2^{k \times n}$ matrix generating an efficiently correctable code}
    \KwIn{$\kpri = (\mathbf{x},\mathbf{y}) \in \mathcal{R}_\omega \times \mathcal{R}_\omega$\newline $(\mathbf{u}, \mathbf{v}) \in \mathcal{R} \times \mathcal{R}$}
    \KwOut{$\ssplit{\mathbf{b}}$, such that $\text{\sc Combine}(\ssplit{\mathbf{b}})=\mathbf{m}\in \mathbb{F}_2^k$}

    $\ssplit{\mathbf{y}} \gets \text{\sc Split}(\mathbf{y})$\;
    $\ssplit{\mathbf{uy}} \gets (\mathbf{u}\cdot\mathbf{y}_0,\ldots,\mathbf{u}\cdot\mathbf{y}_{x-1})$\;
    $\mathbf{b} \sample \mathbb{F}_2^k$\;
    $\ssplit{\mathbf{b}} \sample \text{\sc Split}(\mathbf{b})$\;
    $\ssplit{\mathbf{c}} \gets \ssplit{\mathbf{b}}\mathbf{G}$\; 
    $\mathbf{a} \gets \text{\sc Collapse}(\mathbf{v} ,\ssplit{\mathbf{c}})$\;
    $\overline{\mathbf{b}} \gets \text{\sc Decode}(\text{\sc Collapse}(\mathbf{a},\ssplit{\mathbf{uy}}))$\tcp*[l]{\color{blue}$\overline{\mathbf{b}} = \mathbf{m} \oplus \mathbf{b} $}
    $\mathbf{b}_0 \gets \overline{\mathbf{b}} \oplus\mathbf{b}_0$\tcp*[l]{\color{blue}this sets $\ssplit{\mathbf{b}}=(\mathbf{m} \oplus \mathbf{b} \oplus\mathbf{b}_0,\mathbf{b}_1, \ldots,\mathbf{b}_{x-1} )$}
    \KwRet $\ssplit{\mathbf{b}}$\;
    \end{algorithm}
\end{minipage}
\vspace{-0.4cm}
\end{figure*}

\noindent\textbf{Masked} \textsf{Decrypt}. The decryption function needs to 
protect the private key component $\mathbf{y}$, which is split (line $1$ of 
Algorithm~\ref{alg:dec}), and subsequently multiplied by the public ciphertext
component $\mathbf{u}$, which does not need splitting (Property~\ref{prop:one}).
We now use Property~\ref{prop:two} to protect the decoding procedure (line $6$)
without computing $x$ times the decoding algorithm, which is the most demanding
portion of HQC.PKE \textsf{Decrypt}.
To this end, we denote with a helper procedure, \textsc{Collapse}$(\mathbf{a}, \ssplit{\mathbf{b}})$ the \texttt{xor} sum of $\mathbf{a}$ with all the shares of
$\mathbf{b}$, done in storage order of the shares themselves, as doing a hierarchical
reduction would lower the number of side channel samples to know for an attacker~\cite{ishai2003private,prouff2013masking}.
We protect the input of the \textsc{Decode} procedure, adding a random codeword
to it. This random codeword is obtained drawing a random information word 
$\mathbf{b}$ (line $4$), share splitting it (line $5$) and encoding all the shares, 
obtaining the share-split codeword $\ssplit{\mathbf{c}}$ (line $6$).
We note that a single \textsc{Decode} computation is performed, as opposed
to a straightforward masking scheme which would require $x$ of them to be 
computed.
Thanks to Property~\ref{prop:two}, decoding the sum of share-split codewords
will result in the sum of the corresponding share-split information words.
We therefore collapse the share-split codeword $\ssplit{\mathbf{c}}$ onto the
public decoder input, $\mathbf{v}$, while also performing the subtraction 
$\mathbf{v}-\mathbf{uy}$ (lines $6$,$7$, recalling that the subtraction 
among binary polynomials amounts to a coefficient-wise \texttt{xor}).
The result of the \textsc{Decode} procedure $\mathbf{b}_0$ is thus the message 
information word $\mathbf{m}$, with a random information word $\mathbf{b}$ 
added to it.
Adding such a value to the first share of the share-split representation of
$\ssplit{\mathbf{b}}= (\mathbf{b}_0,\mathbf{b}_1, \ldots,\mathbf{b}_{x-1} )$ turns such a representation into a share split representation of
$\mathbf{m}=(\mathbf{m} \oplus \mathbf{b} \oplus\mathbf{b}_0,\mathbf{b}_1, \ldots,\mathbf{b}_{x-1} )$ which is the value returned by the protected HQC.PKE \textsf{Decrypt}.
Such a value would be recombined, in case the HQC.PKE were used alone; instead, 
considering the HQC.KEM, can be forwarded to the protected components of the FO
construction for further processing.

%
%-------------------------------------------------------------------------------
%
\section{Experimental Evaluation}\label{sec:expres}
In this section we report the results of the efficiency and effectiveness of our
proposed masked implementation.
% In our implementation we will consider only the case $d=2$.
To this end, we choose the Nucleo L$476$RG board, which is endowed with a
STM$32$L$476$RG Cortex-M$4$ microcontroller (which can be clocked up to $80$MHz), 
$128$kiB of RAM, $1$MiB of Flash memory, and a true random number generator.
The choice of the platform was made as it is part of the \texttt{pqm4} benchmarking
project~\cite{PQM4}, allowing us fair comparisons with the current state of the art, 
and the chosen microcontroller is the one picked by NIST as reference platform.
%
The available memory on the microcontroller constrains us to the 
\texttt{hqc-128} parameter set alone and a $x=2$ out of $2$ share splitting, as 
larger ones have operands which do not fit, once share-split, into the RAM of
the microcontroller. 
We employed \texttt{libopencm3} $1$.$0$.$0$ as our hardware abstraction layer 
library and compiled our code with \texttt{gcc-arm-none-eabi} $7$.$2$.$1$
using the \texttt{-O3} optimization flag. 
We performed our evaluation clocking the microcontroller with frequencies in the 
$\{16, 32, 48, 64, 80\}$MHz set to elicit possible timing vulnerabilities due
to different Flash memory wait states.
% We run the same experiments on a x$86$ machine, equipped with an AMD Ryzen $25000$U, $8$GiB of DDR$4$ DRAM and running ArchLinux $6$.$2$; the executables are compiled with
% \texttt{gcc} $12$.$2$.$1$ using the \texttt{-O3} optimization flag and, for AVX implementations (marked with a `$\dagger$'), the \texttt{-march=native} flag. 

\noindent\textbf{Side-channel Security Evaluation.} 
To detect the presence of an information leakage on either the timing or 
the power consumption side channel we employ the well-established Test Vector 
Leakage Assessment (TVLA)~\cite{DBLP:journals/jce/SchneiderM16} procedure,
based on Welch's $t$-test. 
The technique collects two sample populations, one from the cryptographic 
algorithm running on a fixed value as an input, the other running on randomly
generated inputs, while the cryptographic-key in both cases is a fixed value.
TVLA subsequently employs a Student's $t$-test built with the null-hypothesis
that the populations have different means (i.e., some information leakage takes
place). The widely accepted bar for the $t$-test to reject the null-hypothesis
and hence indicate absence of side channel leakage is that the value of the $t$
statistic is in the range $[-4.5;\, 4.5]$.
Small $t$ values give (statistical) evidence to reject the null hypothesis and, 
as a consequence, to be assured that the two populations have the same mean. 
Specifically, values of $t \in [-4.5;\,4.5]$ allow, with a confidence level of 
$99.999\%$, the null hypothesis to be rejected, and then to be highly confident 
that the side-channel behavior of the target implementation fed with random plaintexts 
cannot be distinguished from the one obtained running the target implementation with a 
fixed plaintext. 
\begin{table}[!t]
    \centering
    \caption{$t$-values for the HQC.KEM algorithms. The `$\bullet$' in the \textbf{CT?} column indicates that the multiplication algorithm employed is constant-time \label{tab:ct}}
\resizebox{0.99\textwidth}{!}{    
%     \begin{tabular}{c | c c c | c c | c c | c c | c c }
%         \toprule
%         & & & &  \multicolumn{2}{c|}{\textbf{Encrypt}} & \multicolumn{2}{c|}{\textbf{Decrypt}} &  \multicolumn{2}{c|}{\textbf{Encaps}} & \multicolumn{2}{c}{\textbf{Decaps}}\\
%         \textbf{Round} & \textbf{Version} & \textbf{Mul} & \textbf{CT?} & \texttt{x$86$} & \texttt{ARM} & \texttt{x$86$} & \texttt{ARM} & \texttt{x$86$} & \texttt{ARM} & \texttt{x$86$} & \texttt{ARM} \\
%         \midrule
%         \multirow{2}{*}{3}
%         & \tt{reference} & Shift\&Add              & $\circ$   & $10.65$& $29.19$& $5.59$ & $39.23$& $58.42$& $28.65$& $20.71$& $18.44$\\
%         & \tt{sparsemul} & CT Shift\&Add$^\dagger$ & $\bullet$ & $9.38 $& $23.12$& $7.59$ & $49.56$& $6.75 $& $14.29$& $5.81 $& $23.70$\\
%         \midrule
%         \multirow{3}{*}{4}
%         & \tt{reference} & Karatsuba$^\dagger$     & $\bullet$ & $1.27$ & $1.04$ & $0.97$ & $0.30$ & $3.77$ & $4.01$ & $4.56$ & $0.26$\\
%         & \tt{sparsemul} & CT Shift\&Add$^\dagger$ & $\bullet$ & $0.29$ & $0.07$ & $1.51$ & $1.64$ & $1.58$ & $2.17$ & $1.04$ & $0.22$\\
%         & \tt{fastmul}   & Frobenius FFT           & $\bullet$ & $3.36$ & $0.56$ & $1.70$ & $1.99$ & $4.31$ & $4.25$ & $2.41$ & $0.53$\\
%         \bottomrule
%     \end{tabular}
    \begin{tabular}{c | c c  |  c |  c |  c |  c }
        \toprule
        \textbf{NIST Round} & \textbf{Mult. Alg.} & \textbf{CT?} &  \textbf{Encrypt} & \textbf{Decrypt} &  \textbf{Encaps} & \textbf{Decaps}\\
        \midrule
        \multirow{2}{*}{3}
        & Shift\&Add              & $\circ$   & $29.19$& $39.23$& $28.65$& $18.44$\\
        & CT Shift\&Add           & $\bullet$ & $23.12$& $49.56$& $14.29$& $23.70$\\
        \midrule
        \multirow{3}{*}{4}
        & Karatsuba-based         & $\bullet$ & $1.04$ & $0.30$ & $4.01$ & $0.26$\\
        & CT Shift\&Add           & $\bullet$ & $0.07$ & $1.64$ & $2.17$ & $0.22$\\
        & additive-FFT            & $\bullet$ & $0.56$ & $1.99$ & $4.25$ & $0.53$\\
        \bottomrule
    \end{tabular}    
    }
\vspace{-0.4cm}
\end{table}

To measure the execution times we use the {\em Data Watchpoint and 
Trace Cycle Count Register} (\texttt{DWT\_CYCCNT}), a $32$-bit register that 
counts the number of clock cycles elapsed since its last reset, with a resolution
of one clock cycle; to avoid overflows, we reset the register after all the 
measurements taken from each execution of the target implementation.
Table~\ref{tab:ct} reports the values of the $t$ statistic for the implementations 
of the entire HQC.KEM as provided by the authors to NIST for both Round $3$ and Round $4$ of the competition. For the sake of evaluation and comparison we tested their implementations also replacing our multiplication algorithms based on a constant time shift and add strategy as well as on the additive-FFT. The $t$ values reported in Table~\ref{tab:ct} have been obtained gathering $1,000$ samples for each population (i.e., fixed and randomized plaintexts), and setting the microcontroller frequency to $80$MHz. 
The obtained results are analogous picking any working frequency in the
in the $\{16, 32, 48, 64, 80\}$MHz set, and are omitted for lack of space.
We note that reporting a pass in the $t$-test for the entire HQC.KEM is useful to confirm 
that the entire implementation is constant time. 
Indeed, Guo et. al~\cite{DBLP:journals/tches/GuoHJLNS22} highlighted that
also components besides the multipliers may expose timing information leakage.
This fact is highlighted by both implementations of the Round $3$ HQC proposal
having reported $t$ statistics outside the acceptable interval $[-4.5;4.5]$ (e.g., $23.12$ from the execution of even the {\sc Encrypt} primitive alone, despite the use of a constant time multiplication algorithm).
For the one employing the variable-time shift-and-add multiplication 
the effect is more evident, while the one employing a constant-time multiplication
algorithm has its timing leakage stemming from the sparse polynomial sampler proposed for the Round $3$ HQC.KEM 
implementation, as reported in~\cite{DBLP:journals/tches/GuoHJLNS22}. 
All Round $4$ HQC.KEM implementations, employing
a Karatsuba-based multiplication strategy (as proposed by the HQC authors), or either a constant-time (CT) shift-and-add, or additive-FFT multiplication algorithm (proposed by us) have values for the $t$ statistic within the required interval.


%  As already mentioned we employ the TVLA and take as populations the execution times
% of the algorithm under test, running it with variable and fixed inputs; both populations are sampled $1,000$ times and, for ARM,
% the frequency is set to $80$MHz. 
% The results of the tests are reported in Table~\ref{tab:ctmult}.
% \begin{table}[!t]
%     \centering
%     \caption{$t$-values for the multiplication algorithms\label{tab:ctmult}}
%     \resizebox{0.99\textwidth}{!}{
%     \begin{tabular}{ccccc}
%         \toprule
%         &  {\bf HQC reference impl.} & {\bf HQC optimized impl.}     & \multicolumn{2}{c}{{\bf Our realization}} \\
%         & \textbf{Shift\&Add       } & \textbf{Toom-Cook+Karatsuba} & \textbf{Shift\&Add}  & \textbf{Additive-FFT} \\
%         & {\bf variable-time}    & {\bf constant-time}          & {\bf constant-time}  & {\bf  constant-time}  \\            
%         \midrule
%         \texttt{ARM} & $55.75$ & $0.00$ & $0.00$ & $0.00$ \\
%         \texttt{x86} & $20.94$ & $3.64$ & $3.37$ & $1.23$\\
%         \bottomrule
%     \end{tabular}
%     }
% \end{table}
% 
% As expected, the schoolbook Shift\&Add does not run in constant time while the 
% other algorithms do, as they perform always the same
% operations regardless of the inputs; the constant-time behavior of a part of 
% the system however does not guarantee that the whole
% system runs in constant time. In~\cite{DBLP:journals/tches/GuoHJLNS22}, the 
% authors show that the algorithm employed in HQC and BIKE to sample polynomials
% with a fixed Hamming weight is not constant-time and can be exploited to fully 
% recover the secret key; this vulnerability was discovered
% on the Round $3$ version, and it has been fixed in the subsequent round. 
% We 
% therefore run the TVLA on the PKE and KEM schemes with
% the same settings as before, and find the results reported in Table~\ref{tab:ct}.
\begin{figure}[!t]
\begin{center}
    \subfloat{\includegraphics[scale=0.35]{figures/decode_nomask_50k_t.png}}
    \subfloat{\includegraphics[scale=0.35]{figures/decode_masked_50k_t.png}}
    \newline\noindent
    \subfloat{\includegraphics[scale=0.35]{figures/rsdecode_nomask.png}}
    \subfloat{\includegraphics[scale=0.35]{figures/rsdecode_masked.png}}
\end{center}
    \caption{$t$-value statistics on the electromagnetic traces on the protected RM-RS decoder (top), and zoom on the RS decoder alone (bottom). Figures on the left are from the unprotected implementations, figures on the right from the protected one. All results obtained with $50$k traces.\label{fig:ttest_em}}
\vspace{-0.5cm}
\end{figure}
We employ the electromagnetic emissions as a proxy of the power consumption, 
and measure them with a PicoScope $5244$D oscilloscope sampling at $500$MS/s the 
output of an H-probe placed on top of the STM$32$L$476$RG package, amplifying the
signal with two Agilent INA-$10386$ amplifiers a $52$dB gain.
We collect $50$k samples for each population, and run the $t$ test timewise, 
as per the TVLA~\cite{DBLP:journals/jce/SchneiderM16} methodology, obtaining
time sequences of $t$-statistics. 
Figure~\ref{fig:ttest_em} reports plots of the $t$ statistic as a function of time
for both the unprotected (left) and protected (right) implementations.
In particular, the topmost two figures correspond to the computation of the
RMRS decoder, while the bottom ones are a zoom on the RS decoder alone.
The spike in the value of the $t$ statistic, reaching well beyond the admissible
limit for the unprotected implementation corresponds, in time to the calculation 
of the error corrector polynomial, which was also reported 
by~\cite{schamberger2022power} as an information leaking computation.
By contrast, the time series of the protected implementation are within the $[-4.5;4.5]$ interval, depicted by dashed red lines, showing
the effectiveness of our countermeasure.

% The aim of the test is to provide a quantitative measure of how much the means of two
% populations differ; the t-test statistic is computed as follows:
% \vspace{-10pt}
% \begin{equation}
%     t = \frac{\mu_0 - \mu_1}{\sqrt{\frac{s_0^2}{n_0} + \frac{s_1^2}{n_1}}}
% \end{equation}  
% \vspace{-2pt}
% where $\mu, s^2$ and $n$ are the mean, variance and number of samples of the two populations, respectively.

%%%%%%%%%%%%%%%%%%%%%%%%%%%%%%%%%%%%%%%%%%%%%%%%%%%%%%%%%%%%%%%%%%%%%%%%%%%%%%%%

\begin{figure}[!t]
\input{figures/efficiency.tex}
\caption{Execution times of the encryption, decryption, encapsulation and 
decapsulation algorithms of the \texttt{HQC.KEM} on the \texttt{ARM} 
architecture. The measured components of the \texttt{HQC.KEM} include both the 
non-constant time sampler from Round $3$ (R$3$) and the constant time from 
Round $4$ (R$4$). \texttt{AFFT} marks the additive-FFT multiplier based implementations\label{fig:timing_hqc_kem}}
\vspace{-0.4cm}
\end{figure}
\noindent\textbf{Efficiency evaluation.} We evalutated the computational
efficiency of all the considered multiplication approaches reported in 
Table~\ref{tab:ct}. 
Figure~\ref{fig:timing_hqc_kem} reports the number of 
clock cycles taken by a full execution of the HQC.KEM algorithms, considering
their masked implementation (depicted with patterned columns), 
and provide the results of the unprotected implementation only for the most 
efficient, fully constant time proposal (i.e. the one based on the Round $4$ 
sampler and FFT-based multipliers).
The results in Figure~\ref{fig:timing_hqc_kem} show how the FFT-multiplier based
HQC.KEM is almost an order of magnitude faster than the corresponding Karatsuba
multiplier based one (note that the plot is in logarithmic scale on the $y$ axis).
Our FFT-based multiplier approach also consistently outperforms the 
constant-time shift-and-add one, albeit by a smaller margin.
Finally, we note that the performance overhead which our algorithmic approach
to compute the HQC.KEM as masked is consistently within the $2\times$ range
which is expected from its sub-linear cost in the number of shares (two, in
the figures we report).
Consistently with~\cite{DBLP:journals/iacr/HeinzKLPSS22}, we also report that
the main performance overhead comes from the additional multiplications, the 
impact of the masking on the hash function is negligible with respect to 
the total computation time.


% \newcommand{\ARMtimingplot}[1]{
    \begin{tikzpicture}
        \centering
        \begin{axis}[
            height=70mm, width=150mm,
        %
            bar width=0.1,
            ybar=2pt,
            enlarge x limits={abs=0.5},
            %ymin=0,
            %ymax=170000,
            %ytick distance=10000,
            yticklabel style={font=\small},
            yticklabel=\pgfmathprintnumber{\tick},
            ylabel={Execution time (kcycles)},
            xlabel={Wait States},
            xtick=data,
            xticklabels from table={#1}{Label},
        %
            %nodes near coords,
            %nodes near coords style={font=\tiny,
            %                        /pgf/number format/.cd,
            %                                    precision=0,
            %                                    zerofill,
            %                        },
            legend style={legend pos=north east,
                        cells={anchor=west},
                        font=\footnotesize,
                        }
        ]
        \addplot table [x expr=\coordindex,y=Reference (R3)]{#1};
        \addplot table [x expr=\coordindex,y=sparsemul (R3)]{#1};
        %\addplot table [x expr=\coordindex,y=Reference (R4)]{#1};
        \addplot table [x expr=\coordindex,y=sparsemul (R4)]{#1};
        \addplot table [x expr=\coordindex,y=FAFFT (R4)]{#1};
        \legend{Reference (R3), sparsemul (R3), sparsemul (R4), FAFFT (R4)}    
        \end{axis}
    \end{tikzpicture}
}
% \pgfplotstableread[col sep=comma]{  % enc
    Label, Reference (R3), sparsemul (R3), Reference (R4), sparsemul (R4), FAFFT (R4)
    0,      2726,          18920,           136262,         18965,          8834
    1,      2748,          18937,           138227,         18991,          10107
    2,      2799,          18986,           142093,         19049,          11595
    3,      2886,          19064,           147971,         19154,          13189
    4,      3035,          19196,           155762,         19324,          14856
}\armEncTiming % enc

\pgfplotstableread[col sep=comma]{  % dec
    Label,               Reference (R3), sparsemul (R3), Reference (R4), sparsemul (R3), FAFFT (R4)
    0,                 1942,           9010,           68750,          8915,           5043
    1,                 1950,           9018,           69733,          8926,           5677
    2,                 1965,           9027,           71657,          8941,           6409
    3,                 1989,           9045,           74568,          8958,           7176
    4,                 2024,           9071,           78415,          8998,           7958
}\armDecTiming

\pgfplotstableread[col sep=comma]{  % encaps
    Label,           Reference (R3), sparsemul (R3), Reference (R4), sparsemul (R4), FAFFT (R4)
    0,                 3913,           20105,          137873,         20571,          10441
    1,                 3958,           20148,          139835,         20643,          11763
    2,                 4079,           20261,          143631,         20783,          13333
    3,                 4289,           20461,          149685,         21051,          15087
    4,                 4626,           20787,          157722,         21473,          16994
}\armEncapsTiming

\pgfplotstableread[col sep=comma]{  % decaps
    Label,               Reference (R3), sparsemul (R3), Reference (R4), sparsemul (R4), FAFFT (R4)
    0,                 5923,           29182,          206691,         29557,          15546
    1,                 5974,           29230,          209618,         29634,          17504
    2,                 6110,           29361,          215263,         29795,          19804
    3,                 6341,           29578,          224238,         30084,          22326
    4,                 6724,           29926,          236118,         30540,          25015
}\armDecapsTiming

\pgfplotstableread[col sep=comma]{  % mult
    Label,               Reference (R3), sparsemul (R3), Reference (R4), sparsemul (R4), FAFFT (R4)
    0,                 848,            7897,           66777,          7791,           3062
    1,                 848,            7897,           67747,          7791,           3688
    2,                 851,            7897,           69653,          7791,           4404
    3,                 856,            7897,           72537,          7791,           5146
    4,                 862,            7897,           76345,          7791,           5892
}\armMultTiming
  % contains data for the plots
% \pgfplotstableread[col sep=comma]{  % dec
    Label,               Reference (R3), sparsemul (R3), Reference (R4), sparsemul (R4), FAFFT (R4)
    Encrypt,                148,           604,           157,          667,           1212
    Decrypt,                219,           370,           185,          396,           757
    Encaps,                 216,           667,           220,          727,           1278
    Decaps,                 440,           1041,          407,          1145,          2041
}\mydata
  % contains data for the plots



% % As for the x$86$ machine, we use the \texttt{rdtsc} instruction, that returns the current value of the processor's time stamp counter;
% % to get accurate measurements, we lock the frequency of all cores to $2$.$0$GHz using \texttt{cpupower} and disable
% the Turbo Boost feature.
% In both cases, in order to reduce uncertainty, we run the experiments $1,000$ times and average the results.

% As already mentioned, the multiplication operation is the most time-consuming operation for all primitives, hence the choice of
% the multiplication algorithm is crucial to achieve good performance. 
% The execution times of the various multiplication algorithms are shown in Figure~\ref{fig:multTiming}.
% 
% \begin{figure}[!t]
%     \subfloat{\ARMtimingplotmult{\armMultTiming}}
%     \subfloat{\xtimingplotmult{\xMultTiming}}
%     \centering
%     \newline
%     \timingLegend
%     \caption{Multiplication execution times on ARM (left) and x$86$ (right) in kcycles\label{fig:multTiming}}
% \end{figure}
% 
% As expected, the fastest multiplication algorithm on ARM (after the Shift\&Add) is the one using the Frobenius FFT, followed by the
% constant-time version of Shift\&Add. In the graph we can also see how the execution time for the Frobenius FFT is the only one
% that is significantly impacted by the amount of memory wait states: as explained in Section~\ref{sec:fft} this algorithm heavily
% relies on precomputed tables, requiring a higher amount of memory accesses and thus up to $92$\% more clock cycles; finally, 
% we can see that the Karatsuba algorithm takes the longest amount of time.\\
% Things change dramatically on x$86$: on this architecture we can take advantage of AVX instructions, which allow us to carry out
% operations with $256$-bit vectors in a single clock cycle: this benefits the Karatsuba algorithm, which is the fastest one by a large margin,
% while the Frobenius FFT becomes the slowest one.

% \subsubsection{Constant-time execution}
% As stated by the NIST, the constant-time execution is an important property for a standardization candidate, as it increases
% its security against side-channel attacks: it is hence important to assess the presence of timing leaks in our multiplication implementations
% before using them in the PKE and KEM schemes. As already mentioned we employ the TVLA and take as populations the execution times
% of the algorithm under test, running it with variable and fixed inputs; both populations are sampled $1,000$ times and, for ARM,
% the frequency is set to $80$MHz. 
% % The results of the tests are reported in Table~\ref{tab:ctmult}.
% % \begin{table}[!t]
% %     \centering
% %     \caption{$t$-values for the multiplication algorithms\label{tab:ctmult}}
% %     \resizebox{0.99\textwidth}{!}{
% %     \begin{tabular}{ccccc}
% %         \toprule
% %         &  {\bf HQC reference impl.} & {\bf HQC optimized impl.}     & \multicolumn{2}{c}{{\bf Our realization}} \\
% %         & \textbf{Shift\&Add       } & \textbf{Toom-Cook+Karatsuba} & \textbf{Shift\&Add}  & \textbf{Additive-FFT} \\
% %         & {\bf variable-time}    & {\bf constant-time}          & {\bf constant-time}  & {\bf  constant-time}  \\            
% %         \midrule
% %         \texttt{ARM} & $55.75$ & $0.00$ & $0.00$ & $0.00$ \\
% %         \texttt{x86} & $20.94$ & $3.64$ & $3.37$ & $1.23$\\
% %         \bottomrule
% %     \end{tabular}
% %     }
% % \end{table}
% % 
% As expected, the schoolbook Shift\&Add does not run in constant time while the other algorithms do, as they perform always the same
% operations regardless of the inputs; the constant-time behavior of a part of the system however does not guarantee that the whole
% system runs in constant time. In~\cite{DBLP:journals/tches/GuoHJLNS22}, the authors show that the algorithm employed in HQC and BIKE to sample polynomials
% with a fixed Hamming weight is not constant-time and can be exploited to fully recover the secret key; this vulnerability was discovered
% on the Round $3$ version, and it has been fixed in the subsequent round. We therefore run the TVLA on the PKE and KEM schemes with
% the same settings as before, and find the results reported in Table~\ref{tab:ct}.

% 
% We can see that, even with a constant-time multiplication algorithm, the Round 3 version is not constant-time, while the Round 4 version
% runs in constant time with all the multiplication algorithms. 

% \subsubsection{Power analysis}
% Our power measuring setup employs the PicoScope $5244$D oscilloscope sampling at $500$MS/s, two Agilent INA-$10386$ amplifiers in a cascade
% configuration achieving a $52$dB gain and a H-probe, placed on top of the STM32L476RG package. The EM traces are recorded and analyzed on a 
% computer equipped with an AMD A$8$-$6600$K CPU, $32$GiB of DDR$4$ DRAM running Debian $10$. To ensure that the encryption/decryption process runs
% correctly, we generate the inputs on the same machine, send them to the target via the USART$1$ interface and compare the outputs with the
% expected ones; the executables to generate the inputs and verify the outputs are compiled with \texttt{gcc} ($9$.$3$.$0$) on the same machine,
% using the \texttt{-O3} optimization flag. Due to the length of the traces we focus our attention on the RMRS decoder, which to date is the
% main attack target of power attacks on HQC; this also allows us to capture a higher number of traces ($50,000$).

% Our measurements, reported in Figure~\ref{fig:EMtstat1} and Figure~\ref{fig:EMtstat2}, show a peak in the $t$-value on the unprotected version in correspondance with the calculation of the error corrector polynomial, as
% also reported by~\cite{schamberger2022power}; we can also see how this peak disappears in the masked implementation, showing the effectiveness of our 
% countermeasure.


% \begin{figure}[!t]
%     \centering
% 
%     \caption{$t$-value statistics on the electromagnetic traces on the RS decoder alone\label{fig:EMtstat2}}
% \end{figure}

% \subsection{Comparison with other schemes}
% Similarly to ~\cite{alagic2022status} we can compare the performance of our 
% implementation with the other Round $4$ candidates; 
% it is interesting to note that all three candidates are code-based cryptosystems 
% and that they all have a Cortex-M$4$ implementation.
% In Figure~\ref{fig:time} we can see the different amount of time required by 
% our implementations, BIKE~\cite{DBLP:journals/tches/ChenCK21}
% and McEliece~\cite{DBLP:journals/tches/ChenC21} to perform the different 
% operations; for consistency, we consider only the optimized implementations for 
% Cortex-M$4$ with level $1$ parameters. For HQC, we consider the results
% obtained with zero memory wait states, as the authors of the other 
% implementations did.
 


% The results on the HQC.KEM scheme are impacted mainly by the multiplication algorithm, hence the impact of the
% masking on the hash function~\cite{DBLP:journals/iacr/HeinzKLPSS22} is negligible.
% We can see how the McEliece cryptosystem is the fastest one in terms of encapsulation and decapsulation, but 
% requires $2\cdot 10^{12}$ cycles to generate the keypair. With respect to BIKE, the most promising
% candidate in this round, we can see how HQC outperforms it in terms of key generation and decapsulation,
% not only in the fast implementation but also in the masked one; the performance gap is even more evident
% if we consider the time required for an ephemeral key agreement: HQC requires $30.9$ Mcycles and and $50.8$Mcycles
% respectively for the fast and masked implementation while BIKE, without masking, requires $78.1$ Mcycles.

%
%-------------------------------------------------------------------------------
%
%\section{Analysis of Viable Countermeasures}\label{sec:counters}
%\input{counters.tex}
%
%-------------------------------------------------------------------------------
%
\section{Concluding Remarks}\label{sec:conc}
In this work, we presented the first constant-time and fully masked implementation
of the full HQC.KEM, tailored to Cortex-M$4$ microcontrollers, exploring the possible
constant time multiplication techniques.
We proposed a masking approach to protect the entire HQC.KEM
against power consumption side channel attacks with a cost lower than linear
in the number of shares employed in the masking.
Our results shows the practicality of employing a fully masked
HQC on microcontrollers, hafing performance which are competitive with the other
Round $4$ NIST PQC candidates: indeed our HQC masked decapsulation is
faster than the one for a recent unprotected implementation of BIKE~\cite{DBLP:journals/tches/ChenCK21} (full comparison with Round $4$ candidates in the Appendix at   Figure~\ref{fig:other_r4_cand}).

% of HQC offer a good trade-off in terms of memory consumption and execution time, being faster than the (unmasked) BIKE implementation
% and thus making them suitable for standardization.
% In Figure~\ref{fig:time} we can see the different amount of time required by 
% our implementations, BIKE~\cite{DBLP:journals/tches/ChenCK21}
% and McEliece~\cite{DBLP:journals/tches/ChenC21} to perform the different 
% operations; for consistency, we consider only the optimized implementations for 
% Cortex-M$4$ with level $1$ parameters. For HQC, we consider the results
% obtained with zero memory wait states, as the authors of the other 
% implementations did.
% The results on the HQC.KEM scheme are impacted mainly by the multiplication algorithm, hence the impact of the
% masking on the hash function~\cite{DBLP:journals/iacr/HeinzKLPSS22} is negligible.
% We can see how the McEliece cryptosystem is the fastest one in terms of encapsulation and decapsulation, but 
% requires $2\cdot 10^{12}$ cycles to generate the keypair. With respect to BIKE, the most promising
% candidate in this round, we can see how HQC outperforms it in terms of key generation and decapsulation,
% not only in the fast implementation but also in the masked one; the performance gap is even more evident
% if we consider the time required for an ephemeral key agreement: HQC requires $30.9$ Mcycles and and $50.8$Mcycles
% respectively for the fast and masked implementation while BIKE, without masking, requires $78.1$ Mcycles.


% In this paper we provided the first constant-time implementation of HQC and the
% first masked implementation of HQC's public key encryption scheme.


%
%-------------------------------------------------------------------------------
% ---- Bibliography ----
%
% BibTeX users should specify bibliography style 'splncs04'.
% References will then be sorted and formatted in the correct style.

%\clearpage
\bibliographystyle{splncs04}
\bibliography{refs.bib}
%
%
\end{document}
%
%
%%%%%%%%%%%%%%%%%%%%%%%%%%%%%%%%%%%%%%%%%%%%%%%%%%%%%%%%%%%%%%%%%%%%%%%%%%%%%%%

